\begin{prob}\textbf{Actual Midterm 2, Question 3.}\par
\idea{The two shifted kernels give the available eigenvectors. A quadratic product tests whether one factor per eigenvalue is enough; a cubic zero product supplies the upper bound.}
\begin{steps}
\ii The triangular diagonal gives $2,-1$. For $A-2I$, the equations are $y=0$, $-3z=0$, giving $E(2,T)=\mathbb Re_1$. For $A+I$, they are $3x+y=0$, $3y=0$, giving $E(-1,T)=\mathbb Re_3$.
\ii $m_A$ must contain both root factors. Multiplication gives $(A-2I)(A+I)=\left(\begin{smallmatrix}0&3&0\\0&0&0\\0&0&0\end{smallmatrix}\right)\ne0$, but $(A-2I)^2(A+I)=0$. Thus $m_A=(x-2)^2(x+1)=x^3-3x^2+4$; no proper divisor containing both roots works.
\ii The minimal polynomial has a repeated root, so $T$ is not diagonalizable. Independently, its eigenspaces have dimensions $1+1<3$.
\end{steps}\end{prob}
\begin{prob}\textbf{Practice Midterm 2, Question 3(a,b).}\par
\idea{The range of $T-\lambda I$ is smaller and contains every difference $Tv-\lambda v$, so it controls the extra columns when a basis is extended.}
\begin{steps}
\ii In dimension $1$, for any basis vector $u$ there is a scalar $a$ with $Tu=au$, giving a $1\times1$ triangular matrix. The empty basis handles dimension $0$ when it occurs inside induction.
\ii For $n>1$, choose an eigenvalue $\lambda$. The operator $T-\lambda I$ has a nonzero kernel; rank--nullity gives $\dim\im(T-\lambda I)<n$.
\ii If $u=(T-\lambda I)v$, then $Tu=T(T-\lambda I)v=(T-\lambda I)Tv\in U$. This proves invariance.
\end{steps}\end{prob}
\begin{prob}\textbf{Practice Midterm 2, Question 3(c,d).}\par
\idea{A triangular basis inside $U$ already handles its columns; the remaining columns need only lie in $U$ plus their own line.}
\begin{steps}
\ii The restriction $T|_U$ is an operator on the smaller complex space $U$. Induction gives a basis $u_j$ with $Tu_j\in\operatorname{span}(u_1,\ldots,u_j)$; if $U=0$, use its empty basis.
\ii Extend it by $v_1,\ldots,v_r$. For each $j$, $Tv_j=(T-\lambda I)v_j+\lambda v_j\in U+\mathbb Cv_j$, contained in the span through that vector in the extended list.
\ii Every basis image therefore belongs to its initial span. Proposition 59 makes the matrix upper triangular.
\end{steps}\end{prob}
\begin{prob}\textbf{Practice Midterm 2, Question 4.}\par
\idea{Rearrange the constant term of the minimal relation. If that term is zero, cancellation by an assumed inverse would contradict minimal degree.}
\begin{steps}
\ii Write $m_T=\sum_{j=0}^d a_jx^j$, $a_d=1$. If $a_0\ne0$, then $T[-a_0^{-1}\sum_{j=1}^d a_jT^{j-1}]=I$. The bracket commutes with $T$, so it is a two-sided inverse.
\ii If $a_0=0$, write $m_T=xq$. Invertibility would make $Tq(T)=0$ imply $q(T)=0$, contradicting the smaller degree of nonzero $q$.
\ii $R^2=-I$ by substitution. The map is nonscalar, so there is no degree-one relation and $m_R=x^2+1$.
\ii From $R^2+I=0$, $R(-R)=(-R)R=I$. Hence $R^{-1}=-R$, namely $R^{-1}(x,y)=(y,-x)$, checked by both compositions.
\end{steps}\end{prob}
\begin{prob}\textbf{Midterm 2 Review, Problem 4(1).}\par
\idea{Differentiation couples each sine/cosine pair. Real eigenvalue equations on these pairs have no nonzero solution; only constants remain.}
\begin{steps}
\ii The listed functions are independent: their derivative-value or trigonometric coefficient equations separate the constant and the two sine/cosine pairs. More explicitly, they are eigenvectors of $D^2$ with values $0,-1,-4$ grouped by frequency; $1$ is nonzero and each pair $\sin jx,\cos jx$ is independent by evaluating at $0$ and $\pi/(2j)$. The distinct-eigenspace direct-sum argument from Week 6 therefore gives independence.
\ii Write $f=a+b\sin x+c\cos x+d\sin2x+e\cos2x$. The equation $Df=\lambda f$ gives $\lambda a=0$, $-c=\lambda b$, $b=\lambda c$, $-2e=\lambda d$, $2d=\lambda e$.
\ii For real $\lambda$, these imply $(\lambda^2+1)b=(\lambda^2+1)c=0$ and $(\lambda^2+4)d=(\lambda^2+4)e=0$. Thus $b=c=d=e=0$. A nonzero $f$ requires $\lambda=0$ and $a\ne0$, so the only real eigenvalue is $0$, with eigenspace the constants.
\ii Its eigenspace dimension is $1<5$, so it is not diagonalizable over $\mathbb R$.
\end{steps}\end{prob}
\begin{prob}\textbf{Midterm 2 Review, Problem 4(2,3).}\par
\idea{The rotation's scalar equation is $\lambda^2=-1$. A triangular basis would put such a scalar on its diagonal.}
\begin{steps}
\ii $T(x,y)=(y,-x)$ satisfies $T^2=-I$. For $\lambda=i$, $(1,i)$ satisfies $T(1,i)=i(1,i)$; for $\lambda=-i$, $(1,-i)$ satisfies $T(1,-i)=-i(1,-i)$.
\ii The two distinct eigenvalues give independent vectors, so $((1,i),(1,-i))$ is a complex eigenbasis and its matrix is $\operatorname{diag}(i,-i)$, in particular upper triangular.
\ii Over $\mathbb R$, no scalar solves $\lambda^2=-1$, so no real eigenvalue exists. Any nonempty real upper-triangular matrix has real diagonal eigenvalues by Proposition 60; hence no such basis exists.
\end{steps}\end{prob}
